% Anonymous EJPS manuscript. No author block or identifying PDF metadata.
% AI-use declaration is intentionally not inserted until author approval.
% Synchronized with manuscript/full_manuscript.md.
\documentclass[pdflatex,sn-apa]{sn-jnl}
\usepackage[utf8]{inputenc}
\usepackage[T1]{fontenc}
\usepackage{amsmath,amssymb,amsfonts,mathrsfs}
\usepackage{url}
\hypersetup{pdftitle={Machine-Checked Proofs and Formalization Interfaces in Mathematical Physics: The Epistemology of Osterwalder--Schrader Reconstruction},pdfauthor={},pdfsubject={Anonymous manuscript for double-blind review},pdfkeywords={machine-checked proof, mathematical physics, formalization, scientific representation, Osterwalder--Schrader reconstruction, epistemology}}
\setcounter{secnumdepth}{3}
\raggedbottom
\begin{document}
\title{Machine-Checked Proofs and Formalization Interfaces in Mathematical Physics: The Epistemology of Osterwalder--Schrader Reconstruction}
\abstract{Machine-checked proofs establish consequences relative to the definitions and assumptions of a formal system, but scientific reasoning may require results to pass between different formal representations of a physical theory. This paper asks when such a transfer is warranted, using three Lean developments in the Osterwalder--Schrader (OS) programme. OSforGFF constructs a massive Gaussian free field in a Glimm--Jaffe probability-measure formulation; OSreconstruction proves conditional distributional reconstruction theorems from zero-diagonal Schwinger data; and the OSreconstruction Comparator provides a formally checked correspondence between audit and production records. This is a limited positive case of record-level interface warrant: the correspondence preserves the Schwinger-function family and relevant growth assumptions, while the two records share their underlying mathematical definitions. By contrast, the public developments considered here do not establish a complete connection from the particular OSforGFF construction to the input required by OSreconstruction. The analysis distinguishes local proof warrant, source-formalism correspondence, record-level interface warrant, model instantiation, and target-directed inference. It uses \emph{inter-formalization warrant} as a bounded analytical label for the conditions under which a result may be reused across formal representations, not as a new universal theory or proof-transfer protocol.}
\keywords{machine-checked proof, mathematical physics, formalization, scientific representation, Osterwalder–Schrader reconstruction, epistemology}
\maketitle

\section{Introduction}

The increasing use of interactive theorem provers in mathematical physics raises a question about the relation between formal proof and scientific knowledge. A proof assistant such as Lean can verify that a proposition follows from specified definitions and assumptions in a formal system. When a proof is accepted by the kernel, the resulting warrant is deductive: relative to the soundness of the system and its declared trust base, the proposition follows from the assumptions as they are represented in the formal language. This is a substantial achievement, particularly in areas of mathematical physics where arguments depend on long chains of functional-analytic estimates, distributional constructions, and transformations between different mathematical descriptions. Yet the scientific use of a checked theorem may involve several further steps. The formal statement must correspond to the relevant mathematical content of the physical theory; a concrete model must satisfy the theorem's hypotheses; and, where the result is used to make a claim about a physical target, the model-to-target inference must be justified. These relations become especially important when a theory is formalized in more than one mathematical representation.

The present paper investigates this problem through the Osterwalder--Schrader (OS) reconstruction programme. Its central question is: under what conditions can a machine-checked theorem in one formalization support an inference in another formalization of the same physical theory? The question concerns neither the reliability of Lean's kernel in isolation nor the general relation between mathematical models and reality. It concerns the interfaces through which deductive warrant is transferred between formal representations of scientific mathematics. We use the term \emph{inter-formalization warrant} for the warrant supplied by an adequately established correspondence between such representations. The term is intended as an analytical distinction, not as the name of a new general theory of scientific knowledge.

The problem is related to, but not exhausted by, several existing discussions. Work on formalizing chemical physics has shown how Lean can make assumptions in scientific derivations explicit and how common structures can facilitate the reuse of results across theories. Bobbin et al. (2024), for example, develop reusable structures for gas-law proofs and formalize kinematic equations using calculus. They also distinguish mathematical verification from the empirical adequacy of the theories being formalized. Their work establishes the practical importance of reusable formal structures; the present paper asks a different question about the warrant carried by correspondences between formal records and theorem interfaces, including what must be established when the records originate in different developments. At a more abstract level, the formal-logical literature already provides a precise account of proof-preserving translations. Rabe (2017) defines theory morphisms as compositional translations between formal theories and states judgment preservation in his framework: a derivable source judgment is carried to a derivable target judgment. A morphism from a theory Θ to a theory Θ′ also induces a reduct taking models of Θ′ to models of Θ. This establishes the formal principle that theorem transfer can be justified by a proof-preserving interpretation. The present paper does not propose a new theory morphism or a new proof-transfer mechanism; it asks what additional epistemic work is required when the theories being related are formal representations of a physical formalism and the resulting mathematical claims are used in scientific reasoning.

Their concern is adjacent to a technical literature on proof transfer and heterogeneous formal reasoning. Al Wardani et al. (2023) develop the Distributed Assertion Management Framework, in which formal assertions can be exchanged with explicit provenance and configurable trust in their creators and tools. Guilloud et al. (2025) use SC-TPTP to exchange first-order derivations with equality and demonstrate reconstruction of supported prover outputs in Lean 4. These projects establish that proof exchange, translation, and trust management are already active technical research areas. The present paper does not claim to introduce proof transfer or interoperability. It examines a different epistemic question: when the formal records being connected encode distinct mathematical presentations of a physical theory, what must be established about their mathematical and physical content before a checked result can support scientific inference?

This issue is particularly visible in recent work on the formalization of quantum field theory. Douglas, Hoback, et al. (2026) explain that the same physical ideas can be expressed in mathematically different ways and that formalization requires concrete choices influenced by the definitions and theorems already available in Lean libraries. Their formalization of the massive free field follows the Glimm--Jaffe probability-measure formulation of the OS axioms, rather than the original Schwinger-function formulation. They explicitly note that this is a variant of the original OS axioms and explain the relation between their use of Schwartz functions and the compactly supported test functions in the Glimm--Jaffe presentation. Their paper identifies OS reconstruction as a subsequent milestone for the four-dimensional free-field formalization. The pinned OSforGFF library (Cherkis et al., 2026) separately supplies a corollary, gaussianFreeField\_satisfies\_all\_OS\_axioms\_of\_dim, using the canonical proper-time propagator for every D ≥ 2; separate results give closed-form covariances in D = 2, 3, 4, and 5. Its mathematical representation remains the probability-measure and generating-functional formulation.

Related work on autoformalization has identified a further source of difficulty. Meadows et al. (2026) characterize forms of semantic drift that occur when informal scientific statements are translated into Lean, including notational collapse and abstraction elevation. Their focus is the human-in-the-loop autoformalization of physics problems, whereas the present study concerns established mathematical-physics developments and their relations to one another. Nevertheless, their analysis reinforces the point that kernel correctness is a property of the formal statement actually encoded; it does not alone establish that the statement preserves the intended scientific content. The closest philosophical antecedent is DeDeo and Duede's (2026) analysis of correspondence between mathematical proofs and formal derivations. They distinguish adequate representation of the original theorem from tracking the inferential steps of the informal proof. Their discussion goes further than a simple translation problem: their account of library lock-in and shared models shows how inherited formal infrastructure can hybridize fragments from distinct modelling traditions, while their analysis of structural, causal, and explanatory routes identifies different kinds of evidence for correspondence. The present paper therefore does not claim to have discovered that library history shapes formal representation, or that machine checking alone settles semantic adequacy. It changes the unit of analysis from an informal proof and its formal derivation to two formal representations of mathematical-physical data. The question is whether objects, predicates, test-function domains, quantitative hypotheses, and output witnesses are related in a way that permits a theorem in the source representation to be used in the receiving one. This is a narrower, complementary interface problem; it does not replace their account of proof correspondence.

The philosophical background is provided by work on scientific representation and mathematical practice. Hughes's (1997) account of denotation, demonstration, and interpretation (DDI) distinguishes the representational relation between a model and its target from the reasoning performed within the model and the interpretation of the resulting conclusions. Gelfert (2011) identifies limitations of both DDI and matching-model accounts and argues that mature mathematical formalisms provide heterogeneous inferential and interpretative resources. Gelfert (2022) develops the epistemic role of notation and its rules in mathematical practice. This point is important here: machine-checked proof should not be identified with scientific demonstration as a whole. It is one formalized component of reasoning conducted within a mathematical formalism. Gelfert's account of rigorous results and cross-model warrant transfer (2009) provides a further antecedent. It shows how mathematically established relations can transmit empirical warrant between physical models. The present paper does not claim that such transfer is new; it asks how formal proof and the interfaces between formalizations contribute to the deductive part of that process.

The distinction between local proof warrant and inter-formalization warrant is also consistent with the epistemological heterogeneity of computational science discussed by Imbert and Ardourel (2023). Verification can serve different epistemic roles depending on the object and purpose of the computational practice. A machine-checked proof may warrant a formal consequence, while correspondence with a source theory, instantiation in a model, and empirical assessment remain distinct tasks. The OS case makes these tasks unusually explicit because the projects expose their formal records, theorem hypotheses, test-function spaces, and proof-checking interfaces.

Three public Lean developments are examined. OSforGFF (Cherkis et al., 2026) constructs the massive Gaussian free field as a probability measure on tempered distributions and proves a dimension-generic package of Glimm--Jaffe OS properties. OSreconstruction (Yin et al., 2026) formalizes forward and reverse reconstruction results between Euclidean Schwinger data on a zero-diagonal Schwartz test-function space and Wightman-function distributions on the full Schwartz space. The OSreconstruction Comparator (Douglas, Yin, et al., 2026) provides a formally checked correspondence between audit and production records within the reconstruction development.

The case yields two contrasting findings at different interface levels. First, the Comparator provides a limited positive example of record-level interface warrant: explicit mappings preserve the Schwinger-function family, relate the relevant growth assumptions, and connect the forward, OS-II, and reverse reconstruction statements. Because the audit and production records share their underlying mathematical definitions, this is a checked interface within a common development, not a bridge between independently authored physical formalisms or a second derivation. Second, the public developments considered here do not establish an end-to-end construction connecting the particular OSforGFF probability measure to the OS and growth-condition records required by OSreconstruction. The two projects use different mathematical records and test-function domains. In particular, OSforGFF's OS1 regularity condition is not the arity-dependent OS-II growth condition required by OSreconstruction. Although the OSreconstruction authors report massive-free-field and Proca-Wick-square examples in a private companion project, that report does not identify the massive-free-field example with the specific OSforGFF measure considered here.

The paper's contribution is a source-critical analysis of the conditions under which results can be reused across formal representations in mathematical physics. It distinguishes (i) the local deductive warrant supplied by kernel checking, (ii) correspondence between a source physical formalism and its formal encoding, (iii) record-level interface warrant established by proved mappings between formal records, (iv) instantiation of theorem hypotheses by a concrete model, and (v) target-directed inference from a model to a physical system. These are not stages that every scientific claim must traverse in a single linear sequence. They are distinct relations that must be assessed whenever a claim relies on the composition of formalized results. The analysis does not propose a new general theory of scientific representation, proof correspondence, or empirical warrant transfer. Rather, it isolates the interface obligations that arise when machine-checked results are reused across distinct formal representations of a physical theory.

The method is a conceptual and source-critical case study of the mathematical results and publicly available formalization records in the OS projects, read alongside the original OS literature. It does not present a new formalization or numerical calculation. The mathematical results are attributed to their authors; the paper's contribution lies in the epistemological analysis of the conditions under which these results can be composed. Section 2 reconstructs the relevant mathematical interfaces in the OS projects. Section 3 develops the distinction between local proof warrant and inter-formalization warrant and relates it to scientific representation and model-based inference. Section 4 summarizes the implications for the epistemology of formalized mathematical physics.

\section{The Osterwalder--Schrader Formalization Network}

\subsection{The mathematical setting}

The Osterwalder--Schrader programme establishes conditions under which Euclidean field-theoretic data can be related to relativistic quantum field theory. Its Euclidean objects are Schwinger functions; the reconstructed relativistic objects are Wightman functions or distributions. The original OS papers formulate axioms for Euclidean Green functions and develop the corresponding reconstruction theorem (Osterwalder \& Schrader, 1973, 1975). Osterwalder and Schrader (1975, pp. 282, 287--288) identify a gap in the 1973 forward proof and leave open there whether E0--E4 alone suffice. Their Section IV reconstruction theorem assumes E0′ together with E1--E4 and yields the output bound R0′.

The mathematical setting is not unique. Glimm and Jaffe (1987) formulate the Euclidean theory in terms of a probability measure on field configurations and its generating functional. The original OS formulation, by contrast, is expressed in terms of Schwinger functions as distributions on an appropriate test-function space. Douglas, Hoback, et al. (2026) choose the Glimm--Jaffe formulation for their Lean formalization because it makes the construction of a Euclidean field theory a problem about a functional measure, a mathematical structure well supported by Mathlib. They explicitly note that this is a variant of the original OS axioms and explain the relationship between their use of Schwartz functions and the compactly supported test functions in the Glimm--Jaffe presentation. Glimm and Jaffe (1987, p. 91, n. 1) also discuss the Schwartz-space extension and the free-field measure on tempered distributions. The three formalization projects considered below should therefore be understood as related representations of the OS programme, not as interchangeable encodings of a single record.

\subsection{OSforGFF: the probability-measure formulation}

OSforGFF (Cherkis et al., 2026) constructs the massive Gaussian free field as a probability measure on the space of tempered distributions \(S'(\mathbb R^D)\), where \(D\geq2\) is the spacetime dimension. The general result applies given a propagator, while a further result supplies the canonical proper-time propagator for every \(D\geq2\). Separate results also give closed-form covariances in dimensions \(D=2,3,4,5\). The constructed measure satisfies the project's Glimm--Jaffe formulation of the OS properties.

The axioms are predicates on the probability measure and its generating functional. OS0 states analyticity of the generating functional in finite collections of complex test-function parameters. OS1 is a regularity condition: there exist \(p\in[1,2]\) and \(c>0\) such that

\[
|Z_\mu(f)| \leq
\exp\left[
c\left(
\int_{\mathbb R^D}|f(x)|\,dx+
\int_{\mathbb R^D}|f(x)|^p\,dx
\right)
\right],
\]

together with a subsidiary local-integrability condition on the two-point function when \(p=2\). For the massive Gaussian free field, the proof uses \(p=2\) and \(c=1/(2m^2)\), and establishes the required two-point condition by identifying the mollified two-point function with the covariance kernel away from coincidence. OS2 expresses invariance under the Euclidean group. OS3 is reflection positivity in a complex generating-functional formulation, with time reflection and complex conjugation. OS4 is represented through clustering and ergodicity; the latter is initially formulated for a specified class of finite exponential functionals and established using a polynomial-clustering estimate.

These choices are not merely notational. The OSforGFF documentation and the accompanying paper explain how the formalized conditions relate to the Glimm--Jaffe presentation. In particular, the use of Schwartz test functions and tempered distributions is stated to imply the corresponding compactly supported test-function formulation, since \(D(\mathbb R^D)\subset S(\mathbb R^D)\) and \(S'(\mathbb R^D)\subset D'(\mathbb R^D)\). The formalized ergodicity condition is initially stated for finite exponential functionals. Douglas et al. (2026, Remark 3.2) explain how density in \(L^1(\mu)\), the \(L^2\)-to-\(L^1\) estimate on a probability space, and invariance of \(\mu\) under time translations extend the time-average result to general integrable observables. This extension is attributed to the authors' argument rather than presented as a direct quotation from the accessible pages of Glimm and Jaffe.

The OS1 condition is particularly relevant to the comparison with OSreconstruction. It is the Glimm--Jaffe exponential regularity bound on the generating functional, with an additional two-point integrability condition. It is not the same predicate as the arity-dependent \(E_0'\) growth condition required by the forward OSreconstruction theorem. The former controls the generating functional in terms of \(L^1\) and \(L^p\) norms of a single test function; the latter is a quantitative estimate on the family of \(n\)-point Schwinger distributions, with Schwartz seminorm orders increasing linearly with \(n\). The shared language of regularity or growth does not establish a mathematical implication between them.

\subsection{OSreconstruction: zero-diagonal Euclidean data and conditional reconstruction}

OSreconstruction (Yin et al., 2026) formalizes reconstruction for every positive spatial dimension \(d\), with spacetime dimension \(d+1\). Its Euclidean input is a family of Schwinger functions acting on the zero-diagonal Schwartz test-function space. Its Wightman distributions act on the full Schwartz space. The zero-diagonal domain is mathematically significant: it accommodates the coincidence singularities of Euclidean Schwinger functions without presupposing that they define distributions on the full Schwartz space.

The Euclidean input satisfies temperedness, linearity and reality on the zero-diagonal test space; translation and connected Euclidean rotation invariance; reflection positivity on Borchers sequences supported in the ordered positive-time region; permutation symmetry; and the cluster property. The forward reconstruction theorem additionally assumes a quantitative growth condition. Schematically, this condition takes the form

\[
\|S_n(f)\|
\leq
\alpha\,\beta^n(n!)^\gamma
\|f\|_{ns,ns},
\]

where \(\alpha\) and \(\beta\) are positive, \(\gamma\) is real, and the Schwartz seminorm orders grow linearly with the arity \(n\). The theorem concludes the existence of a Wightman-function family whose distributions are Wick-paired with the input Schwinger family.

An original-coordinate formulation of the OS-II growth condition requires a positive Sobolev index \(s\) and, for \(n>0\), the estimate

\[
\|S_n(f)\|
\leq
\alpha (n!)^\gamma
\|f\|_{\mathrm{OSII},ns},
\]

where the final seminorm is the original-coordinate seminorm on the zero-diagonal test space. A separate normalization is imposed at \(n=0\). A proved conversion takes this condition to the arity-linear formulation, with the parameter choice

\[
\alpha' = \max\{1,\alpha\},\qquad
\beta'=(d+2)^s,\qquad
\gamma'=\gamma+2s.
\]

The stronger OS-II reconstruction result establishes the existence of a Wick-paired Wightman family, the output growth estimate \(R'_0\), and uniqueness among Wightman families satisfying the same Wick-pairing relation. This uniqueness concerns the distribution family, not all possible physical realizations. The quantitative conversion is part of the mathematical result; it is not an identification based on similar terminology.

The reverse theorem constructs Euclidean Schwinger data from a Wightman-function family and proves the corresponding OS properties for that construction. It does not assume the forward growth condition. The two directions therefore do not form an unconditional biconditional between arbitrary OS data and arbitrary Wightman data. The public endpoint is a family of Wightman distributions satisfying the formalized Wightman-function axioms, including covariance, spectral support, locality, positivity, Hermiticity, and clustering. The subsequent GNS construction of a Hilbert space and operator-valued fields lies outside the scope of this formalization.

The public OSreconstruction package also includes a conditional construction of an OS record from a supplied Wightman family. It does not provide a publicly inspectable concrete example that supplies both the forward OS record and the corrected growth condition. The README reports that a private companion project contains examples for the massive free field and the Proca Wick square. These examples are therefore author-reported; the private proof artifacts are not independently inspectable. The report does not identify the massive-free-field example with the particular probability-measure construction considered in OSforGFF.

\subsection{Comparator: a formally checked interface between OS records}

The OSreconstruction Comparator (Douglas, Yin, et al., 2026) provides an audit representation and a corresponding production representation for the forward, OS-II, and reverse reconstruction results. The two representations are checked separately but share the same underlying mathematical definitions. The Comparator therefore offers a useful example of a formally checked interface without constituting an independently authored physical formalization.

Explicit mappings in both directions preserve the Schwinger-function family and carry the relevant growth assumptions between the two records. They also establish the correspondence between the associated Wick-pairing and output-growth conditions. Under these mappings, the forward, OS-II, and reverse reconstruction statements correspond. The interface is thus more than an informal assertion that the two representations mean the same thing: the correspondence itself is part of the formal mathematical result.

The epistemic scope of this result is limited. Because the two representations share their underlying definitions, the check supports reuse between the corresponding records but does not independently establish that those definitions adequately capture the physical OS framework. The Comparator is therefore a positive example of relative, record-level interface warrant, not an independent derivation or an external validation of the reconstruction theory.

\subsection{What is and is not connected in the public formalization network}

The projects exhibit different kinds of correspondence evidence. For OSforGFF, the authors explain how the probability-measure predicates relate to the Glimm--Jaffe axioms, including the relation between Schwartz and compactly supported test functions and the ergodicity argument. For the Comparator, the correspondence between the audit and production records is itself formally established: the mappings preserve the Schwinger-function family and the relevant growth assumptions, and relate the corresponding reconstruction statements. This is a checked interface between records of common mathematical origin, not an independently developed physical formalism.

Across the public developments considered here, no complete construction has been established that starts from the particular OSforGFF probability measure and supplies both the zero-diagonal Schwinger record and the arity-dependent growth condition required by the OSreconstruction forward theorem. Such a construction would require a map from the measure and generating functional to the relevant Schwinger distributions, a proof that the OS predicates transfer to the zero-diagonal formulation, and a witness for the required growth condition. The absence of such a construction in the public materials considered here is not a claim of mathematical impossibility or of its absence from private work.

The significance of the case does not depend on treating the missing public interface as a defect in either mathematical project. Each project establishes the results it states. The issue is whether, and under what conditions, those results can be composed to support a further scientific claim. This is a question about the warrant of the interfaces between formalizations, not about the validity of the local proofs.

\section{Inter-Formalization Warrant and Scientific Knowledge}

\subsection{Local proof warrant and the scope of machine checking}

A machine-checked proof provides a strong form of deductive warrant for a proposition in a specified formal system. If a Lean kernel accepts a proof of

\[
\Gamma\vdash_{\mathcal L}\varphi,
\]

then, relative to the soundness of the formal system and its declared trust base, phi follows from Gamma as represented in L. This warrant is local: it concerns the formal proposition, the definitions through which its objects are represented, the assumptions on which it depends, and the imported declarations and kernel used in checking it. The OS case illustrates that locality does not make the warrant weak. OSforGFF proves a substantial collection of properties for a constructed Gaussian measure; OSreconstruction proves conditional reconstruction theorems; and the Comparator checks the correspondence between a separately elaborated audit surface and the production theorem statements. Each is a significant mathematical achievement.

The locality of proof warrant should not be confused with the claim that formalization is merely syntactic. As Douglas, Hoback, et al. (2026) emphasize, choosing definitions for a physical theory requires mathematical judgment. The available library, the structure of existing theorems, and the tractability of a proof can influence the representation adopted. The resulting formalization may be mathematically appropriate while differing from the most familiar formulation used by physicists. The epistemic task is therefore not to demand a word-for-word translation, but to establish that the formal representation preserves the mathematical content relevant to the claim being made.

DeDeo and Duede's (2026) analysis is relevant not only for its distinction between adequate representation and tracking, but also for its account of library lock-in and shared models. They show how inherited library dependencies can lead formally correct derivations to hybridize fragments from distinct modelling traditions, and they identify structural, causal, and explanatory routes for assessing correspondence. This overlap matters for the present case. A proved record map supplies structural evidence for a relation between formal representations, but it does not by itself settle the provenance of the mapping, its explanatory role, or the adequacy of the source formalism. Their principal relation is between an informal proof and a formal derivation; the relation examined here is between a source formal representation and a receiving formal representation. The question is whether the objects, predicates, test-function domains, quantitative assumptions, and reconstructed witnesses are mapped in a direction and at a level sufficient for theorem reuse. The present analysis is therefore complementary to, and narrower than, their account; it does not claim a new general theory of correspondence or library dependence.

Gelfert's analysis of mature mathematical formalisms helps avoid an overly linear account of this relation. The DDI framework distinguishes denotation, demonstration, and interpretation, but Gelfert (2011) argues that successful scientific representation is often heterogeneous and depends on mature formalisms that supply both inferential and interpretative resources. Machine-checked proof should therefore not be identified with Demonstration as a whole. It is a formally articulated component of reasoning within a mathematical formalism. Nor should Denotation and Interpretation be treated as two final steps that occur only after a proof has been completed. The formalism itself can carry interpretative commitments and locally applicable rules that shape the meaning of the proof's objects and predicates.

\subsection{From formalization correspondence to inter-formalization warrant}

The OS case reveals a further problem when results from distinct formalizations are used together. A proof in one formalization does not automatically become a proof in another, even when both projects concern the same physical theory. The receiving formalization may use different records, test spaces, definitions of symmetry or positivity, or quantitative hypotheses. The warrant of the source theorem can be transferred only when the relevant correspondence between these structures is established.

We use inter-formalization warrant to refer to the warrant supplied by a verified interface through which a result in one formalization is made available in another. The notion is narrower than empirical warrant transfer between physical models and distinct from the general problem of translating an informal proof into a formal derivation. It concerns the relation between two formal representations. At a minimum, an interface must make explicit which objects correspond, how the relevant predicates and hypotheses are translated, and whether the conclusion and any witnesses are preserved. Where the theorem is conditional, the receiving formalization must also establish that its assumptions imply the translated source hypotheses. Where a reconstructed object is produced, the interface must identify the relevant output family rather than merely show that an object of a broadly similar type exists.

This can be expressed schematically. Suppose a source formalization establishes

\[
\Gamma_1\vdash_{\mathcal L_1}\varphi_1.
\]

One sufficient route for a receiving formalization \(\mathcal L_2\) to use the source result for a target proposition \(\varphi_2\) is to provide a translation or interpretation \(\tau\) that preserves the relevant derivation and to prove the required relations between assumptions and conclusions. One schematic audit condition is

\[
\begin{gathered}
\Gamma_1\vdash_{\mathcal L_1}\varphi_1,\qquad
\tau(\Gamma_1)\vdash_{\mathcal L_2}\tau(\varphi_1),\\
\Gamma_2\vdash_{\mathcal L_2}\tau(\Gamma_1),\qquad
\Gamma_2\vdash_{\mathcal L_2}\bigl(\tau(\varphi_1)\rightarrow\varphi_2\bigr)
\\
\Longrightarrow\quad
\Gamma_2\vdash_{\mathcal L_2}\varphi_2.
\end{gathered}
\]

The display is an audit schema, not a universal sufficient-condition theorem for scientific knowledge. The first line represents proof preservation under the translation; the second records hypothesis transport and conclusion preservation in the receiving formalization. The point is not that every interface must be an isomorphism or definitional equality. Depending on the application, it may be an equivalence of records, a structure-preserving map, or a proved implication. What matters is that the direction and strength of the correspondence are explicit and proved at the level required by the intended inference.

The Comparator provides a limited positive example of record-level interface warrant. Its proved mappings preserve the Schwinger-function family and the relevant growth assumptions, allowing the forward, OS-II, and reverse reconstruction statements to be related across the audit and production records. Because these records share their underlying mathematical definitions, the result warrants reuse between them but does not independently validate the physical adequacy of the shared formalization.

The proof-preservation component is familiar from the theory-morphism literature. Rabe (2017), for example, defines theory morphisms and states judgment preservation within his framework; theorem preservation is a special case of that result. The schematic conditions above are therefore not offered as a new formal result. Their role here is diagnostic: they identify the proof-preserving component of an interface so that it can be distinguished from the additional question of whether the translated objects and predicates retain the intended physical interpretation and whether the receiving theorem's hypotheses are instantiated by the relevant model.

This use of the term should be distinguished from technical proof-system interoperability. DAMF records the provenance and dependencies of formal assertions and allows recipients to decide which agents and tools to trust; it does not guarantee the logical compatibility of assertions imported from heterogeneous systems (Al Wardani et al., 2023). SC-TPTP, by contrast, supports the reconstruction and checking in a receiving prover of supported first-order derivations with equality, including demonstrated reconstruction in Lean 4 (Guilloud et al., 2025). These projects therefore address different tasks: provenance-sensitive assertion exchange and proof reconstruction. The present analysis concerns the scientific-epistemic status of interfaces between mathematical records that may encode different formulations of a physical theory. The problems can overlap, but proof reconstruction at the formal level does not by itself establish correspondence between the encoded mathematical objects and the intended physical formalism.

The distinction is also visible in the work of Bobbin et al. (2024). Their reusable structures support gas-law proofs in thermodynamics, while their treatment of kinematics formalizes equations of motion using calculus; these are distinct applications rather than a transfer of gas-law proofs into kinematics. The present analysis does not claim that proof reuse or interoperable scientific structures are new. Rather, the OS case shows why the existence of an interface is an epistemic matter: when a result is reused across representations, the assumptions, predicates, and objects that make the transfer legitimate must be identified and preserved. The issue is especially consequential when the source and receiving formalizations are developed in different projects, use different test-function spaces, or formalize different versions of a familiar axiom system.

\subsection{The OS case and the limits of composition}

The contrast between the Comparator and the OSforGFF-to-OSreconstruction relation should not be read as a pair of symmetric cross-project examples. The Comparator establishes a record-level interface between two representations based on common mathematical definitions. By contrast, the public developments considered here do not provide an explicit construction connecting the particular OSforGFF measure to both the zero-diagonal Schwinger data and the arity-dependent growth condition required by the OSreconstruction forward theorem. This is a statement about the available public formalizations, not a claim of mathematical impossibility. The distinction is not resolved by noting that both developments concern the massive free field or invoke the OS axioms. OSforGFF is formulated through a probability measure and generating functional on Schwartz functions; OSreconstruction takes a family of Schwinger distributions on a zero-diagonal Schwartz test space and requires an arity-dependent growth estimate. The mathematical passage between these records must be supplied.

The difference between the two growth conditions is particularly instructive. OSforGFF proves the Glimm--Jaffe OS1 exponential bound on the generating functional, with \(p=2\) and \(c=1/(2m^2)\) for the massive Gaussian free field, together with the required local integrability of the two-point function. OSreconstruction requires a quantitative estimate on the family of Schwinger distributions whose Schwartz seminorm orders increase with the arity. These are different predicates on different mathematical objects. The first may be relevant to establishing regularity of the generating functional and the existence of moments; the second is an explicit hypothesis of the forward reconstruction theorem. A relation between them may be mathematically provable, but it cannot be inferred from their shared association with regularity or growth.

The OSreconstruction authors report that the corrected growth condition has examples for the massive free field and the Proca Wick square in a private companion project. Since the private proofs are not publicly inspectable, this report cannot establish that the massive-free-field example is the same object as the particular OSforGFF probability measure or supply the missing bridge between the two formalizations.

This analysis also clarifies the relation to Gelfert's account of cross-model warrant transfer. In Gelfert's (2009) example, the limiting relation at half-filling and strong coupling allows evidence for the antiferromagnetic Heisenberg model to support the Hubbard model in the relevant regime. Inter-formalization warrant concerns a different, though related, level: it concerns the proof-preserving correspondence required before a mathematical result in one formal representation can be used in another. Such a correspondence may contribute to a rigorous relation between physical models, but it does not by itself establish the empirical warrant of either model or the physical relevance of the mapping. Machine checking can strengthen the deductive component of a warrant-transfer argument without generating empirical warrant independently.

\subsection{From model-internal knowledge to target-directed inference}

Once a theorem has been transported to the relevant formal representation and its hypotheses have been instantiated by a concrete model, the result supports a model-internal claim. This still does not automatically warrant a claim about an actual physical target. Hughes (1997) distinguishes demonstration in the model from interpretation of its results for a target. Suárez (2004) emphasizes a representation's capacity to license inferences about a target. On Contessa's (2007) account, a valid surrogative inference is sound only if its conclusion is true of that target. Frigg and Nguyen's (2018) DEKI account explains how a key maps properties exemplified by a model to properties that a user imputes to its target. Model adequacy is also purpose-sensitive: a model may be suitable for a mathematical or structural purpose without being adequate for quantitative prediction or empirical description (Parker, 2020).

The OS case should not be taken to imply that every machine-checked result must culminate in an empirical claim. Models may function as epistemic artifacts for investigating possible structures and modal relations, without representing a particular actual system (Knuuttila, 2021). The kind of warrant required depends on the claim being made. A formal theorem, a model-internal result, an explanatory claim, and a target-directed empirical inference are different epistemic achievements. Inter-formalization warrant concerns the transition between formal representations and should not be conflated with the model-to-target relation.

The resulting account is therefore not a linear pipeline in which kernel acceptance is followed automatically by physical interpretation and empirical confirmation. It is a set of distinguishable relations. Kernel checking warrants a derivation in a specified formal system. Formalization correspondence relates the encoded statement to its source mathematical formalism. Inter-formalization warrant relates distinct formal representations and permits theorem reuse when the relevant interfaces are proved. Hypothesis instantiation establishes that a concrete model satisfies the receiving theorem's assumptions. Target-directed inference requires an additional account of representation, purpose, and relevant evidence. These relations may be supported by different kinds of argument, and their requirements depend on the scientific claim at issue.

\subsection{Implications for formalized mathematical physics}

The OS case suggests that the epistemic significance of large formalization projects should not be assessed only by counting checked theorems or examining their local axiom footprints. Such information is important, but it does not describe how results can be reused across formal representations. A mature mathematical formalism often admits several mathematically distinct formulations, and a formalization project necessarily selects one of them. As the project develops, further formalizations may encode related objects and theorems in different records or libraries. The resulting network can support substantial reuse, but its epistemic reliability depends on the quality of the interfaces that connect its components.

This conclusion complements rather than replaces existing analyses. Bobbin et al. (2024) show how common mathematical structures can enable interoperable scientific proofs. Douglas, Hoback, et al. (2026) document how the choice of mathematical formulation is shaped by the available formal infrastructure and by the requirements of the physical theory. Meadows et al. (2026) identify semantic drift in the translation of scientific statements into formal language. DeDeo and Duede (2026) analyze the correspondence between informal proofs and formal derivations, including the effects of library lock-in and shared modelling traditions. The present case connects these discussions at the level of formal-record interfaces: it distinguishes a proved audit-to-production bridge from the separate cross-project instantiation problem and shows what must be established before a checked result can be reused across those interfaces.

The contribution is deliberately limited. It is not a new general theory of scientific representation, a replacement for DDI, a new account of library dependence, or a new account of cross-model empirical warrant. It is a source-critical, formalization-aware analysis of the conditions for reusing machine-checked results across distinct formal records in mathematical physics. The OS network supplies a concrete setting: the Comparator demonstrates a record-level audit-to-production interface, while the OSforGFF-to-OSreconstruction comparison identifies a separate cross-project instantiation obligation that is not established in the inspected public artifacts.

\section{Conclusion}

The Osterwalder--Schrader case shows that the epistemic significance of a machine-checked theorem depends not only on the validity of its proof within a formal system but also on the relations through which the result is connected to other representations of a physical theory. Kernel checking provides local deductive warrant for the formal proposition and assumptions that it verifies. When a result is reused across formalizations, an additional interface must establish the correspondence between the relevant objects, predicates, hypotheses, and conclusions. When the theorem is applied to a concrete model, its assumptions must be instantiated; when the result is used to make a claim about a physical target, the model-to-target inference requires further justification.

The three OS developments illustrate these distinctions. OSforGFF establishes the Glimm--Jaffe probability-measure formulation of the OS properties for every spacetime dimension \(D\geq2\) with the canonical propagator. OSreconstruction proves conditional distributional reconstruction theorems from zero-diagonal Schwinger data, with an arity-dependent growth condition in the forward direction. The Comparator provides a formally checked interface between audit and production records: its mappings preserve the Schwinger-function family and relevant growth assumptions, and relate the forward, OS-II, and reverse reconstruction results. Because the records share their underlying mathematical definitions, this is a record-level result within a common development, not an independent physical formalization. The public developments considered here do not establish an end-to-end connection from the particular OSforGFF measure to the zero-diagonal input and growth condition required by OSreconstruction. This bounded finding does not imply that such a derivation is mathematically impossible.

The analysis therefore distinguishes local proof warrant from inter-formalization warrant. The latter is not a substitute for proof checking; it concerns the mathematical correspondence required for a checked result to be transferred between formal representations. This distinction complements existing work on proof correspondence, interoperable scientific formalizations, mature mathematical formalisms, and cross-model warrant transfer. It also clarifies why the existence of separately verified components does not by itself establish the validity of a composed scientific inference.

\section*{References}
\noindent Al Wardani, F., Chaudhuri, K., \& Miller, D. (2023). Formal reasoning using distributed assertions. In U. Sattler \& M. Suda (Eds.), \emph{Frontiers of Combining Systems (FroCoS 2023)} (Lecture Notes in Computer Science, Vol. 14279, pp. 176--194). Springer. \url{https://doi.org/10.1007/978-3-031-43369-6_10}\par\medskip

\noindent Bobbin, M. P., Sharlin, S., Feyzishendi, P., Dang, A. H., Wraback, C. M., \& Josephson, T. R. (2024). Formalizing chemical physics using the Lean theorem prover. \emph{Digital Discovery}, 3(2), 264--280. \url{https://doi.org/10.1039/D3DD00077J}\par\medskip

\noindent Cherkis, S. A., Douglas, M. R., Hoback, S., Mei, A., \& Nissim, R. (2026). \emph{OSforGFF} [Lean 4 formalization]. GitHub. \url{https://github.com/mrdouglasny/OSforGFF}\par\medskip

\noindent Contessa, G. (2007). Scientific representation, interpretation, and surrogative reasoning. \emph{Philosophy of Science}, 74(1), 48--68. \url{https://doi.org/10.1086/519478}\par\medskip

\noindent DeDeo, S., \& Duede, E. (2026). A correspondence problem for mathematical proof. \emph{Philosophy of Science}, forthcoming. Preprint: arXiv:2603.13680v2.\par\medskip

\noindent Douglas, M. R., Hoback, S., Mei, A., \& Nissim, R. (2026). \emph{Formalization of QFT} [Preprint]. arXiv. \url{https://arxiv.org/abs/2603.15770}\par\medskip

\noindent Douglas, M. R., Yin, X., \& Mei, A. (2026). \emph{OSreconstruction-comparator} [Formal specification and proof-checking project]. GitHub. \url{https://github.com/mrdouglasny/OSreconstruction-comparator}\par\medskip

\noindent Frigg, R., \& Nguyen, J. (2018). The turn of the valve: Representing with material models. \emph{European Journal for Philosophy of Science}, 8(2), 205--224. \url{https://doi.org/10.1007/s13194-017-0182-4}\par\medskip

\noindent Gelfert, A. (2009). Rigorous results, cross-model justification, and the transfer of empirical warrant: The case of many-body models in physics. \emph{Synthese}, 169(3), 497--519. \url{https://doi.org/10.1007/s11229-008-9431-6}\par\medskip

\noindent Gelfert, A. (2011). Mathematical formalisms in scientific practice: From denotation to model-based representation. \emph{Studies in History and Philosophy of Science Part A}, 42(2), 272--286. \url{https://doi.org/10.1016/j.shpsa.2010.11.035}\par\medskip

\noindent Gelfert, A. (2022). Thinking with notations: Epistemic actions and epistemic activities in mathematical practice. In M. Friedman \& K. Krauthausen (Eds.), \emph{Model and Mathematics: From the 19th to the 21st Century} (pp. 333--362). Birkhäuser. \url{https://doi.org/10.1007/978-3-030-97833-4_11}\par\medskip

\noindent Glimm, J., \& Jaffe, A. (1987). \emph{Quantum physics: A functional integral point of view} (2nd ed.). Springer. \url{https://doi.org/10.1007/978-1-4612-4728-9}\par\medskip

\noindent Guilloud, S., Cailler, J., Gambhir, S., Poiroux, A., Herklotz, Y., Bourgeat, T., \& Kunčak, V. (2025). Interoperability of proof systems with SC-TPTP. In C. Barrett \& U. Waldmann (Eds.), \emph{Automated Deduction -- CADE 30} (Lecture Notes in Computer Science, Vol. 15943, pp. 325--340). Springer. \url{https://doi.org/10.1007/978-3-031-99984-0_18}\par\medskip

\noindent Hughes, R. I. G. (1997). Models and representation. \emph{Philosophy of Science}, 64, S325--S336. \url{https://doi.org/10.1086/392611}\par\medskip

\noindent Imbert, C., \& Ardourel, V. (2023). Formal verification, scientific code, and the epistemological heterogeneity of computational science. \emph{Philosophy of Science}, 90(2), 376--394. \url{https://doi.org/10.1017/psa.2022.78}\par\medskip

\noindent Knuuttila, T. (2021). Epistemic artifacts and the modal dimension of modeling. \emph{European Journal for Philosophy of Science}, 11, 65. \url{https://doi.org/10.1007/s13194-021-00374-5}\par\medskip

\noindent Meadows, J., Zhang, L., \& Freitas, A. (2026). FormalScience: Scalable human-in-the-loop autoformalisation of science with agentic code generation in Lean. In \emph{Proceedings of the 64th Annual Meeting of the Association for Computational Linguistics} (pp. 23059--23077). \url{https://doi.org/10.18653/v1/2026.acl-long.1057}\par\medskip

\noindent Osterwalder, K., \& Schrader, R. (1973). Axioms for Euclidean Green's functions. \emph{Communications in Mathematical Physics}, 31(2), 83--112. \url{https://doi.org/10.1007/BF01645738}\par\medskip

\noindent Osterwalder, K., \& Schrader, R. (1975). Axioms for Euclidean Green's functions II. \emph{Communications in Mathematical Physics}, 42(3), 281--305. \url{https://doi.org/10.1007/BF01608978}\par\medskip

\noindent Parker, W. S. (2020). Model evaluation: An adequacy-for-purpose view. \emph{Philosophy of Science}, 87(3), 457--477. \url{https://doi.org/10.1086/708691}\par\medskip

\noindent Rabe, F. (2017). Morphism axioms. \emph{Theoretical Computer Science}, 691, 55--80. \url{https://doi.org/10.1016/j.tcs.2017.07.006}\par\medskip

\noindent Suárez, M. (2004). An inferential conception of scientific representation. \emph{Philosophy of Science}, 71(5), 767--779. \url{https://doi.org/10.1086/421415}\par\medskip

\noindent Yin, X., Douglas, M. R., \& Mei, A. (2026). \emph{OSreconstruction in Lean} [Lean 4 formalization]. GitHub. \url{https://github.com/xiyin137/OSreconstruction}\par\medskip
\end{document}
